\section{Approach}
\label{sec:approach}

\subsection{Hidden-State Extraction}

All probes read the residual stream of Qwen3-8B, run in bf16 precision with
thinking mode explicitly disabled. For a given prompt, we apply the model's
chat template, run a forward pass with \texttt{output\_hidden\_states=True}
and no key-value caching, and extract the hidden state at the last token of
the templated input (before any generation) at every one of the model's 37
residual-stream checkpoints (layer 0, the embedding output, through layer
36, the final block's output). No text is generated at any point in this
extraction; the only computation performed is the forward pass itself.

\subsection{Ambiguity-Degree Probe}
\label{sec:approach:k-probe}

The first probe answers a coarser question: given a coding prompt, how much
unresolved specification ambiguity does it contain, expressed as an ordinal
count $K \in \{1, 2, 3\}$ of independently ambiguous details? We extract the
last-token hidden state of the raw, unmodified prompt at each layer and fit
an independent multinomial logistic regression per layer on top of a
100-component PCA plus standardization pipeline, with the regularization
strength $C$ swept on a validation split. Layer and configuration selection
use grouped cross-validation, splitting by source-problem identifier
(\texttt{StratifiedGroupKFold}) so that no two rows derived from the same
underlying problem appear on both sides of a split -- a leakage guard that
is load-bearing for this task, since near-duplicate prompts sharing a
source problem would otherwise let a probe memorize surface features of a
specific problem rather than learning a general ambiguity signal. We report
ordinal correlation (Spearman's $\rho$ between the predicted expected value
$E[K]$ under the probe's class-probability output and the true $K$) as the
primary metric, alongside majority-class, random, and prompt-length-only
baselines computed on the identical split.

\subsection{Clarifying-Question Ranking Probe}
\label{sec:approach:rank-probe}

The second, and practically more important, probe addresses a different
problem: given a pool of $N$ candidate clarifying questions (with
LLM-elicited answers) for an ambiguous prompt, which one should the agent
actually ask? We treat this as pairwise ranking cast as classification, in
the style of RankSVM/RankNet: each candidate's answer is concatenated onto
the base prompt (\texttt{"\{prompt\}\textbackslash n\textbackslash
nClarification: \{answer\}"}), and a batched forward pass over all $N$
concatenated variants (one GPU call, $N$ sequences) extracts $N$ hidden
states at a single selected layer. A classifier trained on
tier-difference vectors between higher- and lower-quality candidates --
where candidate quality tiers (0/1/2, unrelated/plausible/matches-the-
labeled-ambiguity) are assigned by an independent LLM judge -- scores each
candidate, and the pipeline asks the top-scored one. We additionally
evaluate a supervised-contrastive (SupCon) variant of this probe, which
trains directly on raw (non-PCA-reduced) activations with a contrastive
objective and, as reported in Section~\ref{sec:results}, improves
cross-source generalization at some cost to in-distribution accuracy
relative to the PCA-based classifier.

Critically, this ranking probe's inference cost is $O(N)$ forward passes
executed as a single batched GPU call, not the $O(1)$ single forward pass
of the ambiguity-degree probe in Section~\ref{sec:approach:k-probe} --
each candidate's answer text must be present in the input for its
hidden state to reflect that candidate specifically. The relevant cost
comparison is therefore not ``one forward pass vs.\ one generation call''
but ``one batched forward pass over $N$ short sequences vs.\ one
generation call that reasons over all $N$ candidates inline,'' which we
measure directly in Section~\ref{sec:results:cost}.

\subsection{The Two-Round Elicit-Rank-Clarify-Code Loop}
\label{sec:approach:loop}

We deploy the ranking probe in an agentic pipeline that mirrors how a
clarify-then-code assistant would operate in practice:

\begin{enumerate}
\item \textbf{Elicit}: given an ambiguous prompt, generate a pool of
candidate clarifying questions and, for each, an oracle answer (an
independent LLM answers the question as if it were the ground-truth
requirement author).
\item \textbf{Rank}: score all candidates with the trained probe (or, for
the zero-shot baseline in Section~\ref{sec:experimental_setup}, ask the
coding model to pick directly from a single inlined list) and select the
top-ranked candidate.
\item \textbf{Clarify}: append the selected candidate's question-answer
pair to the prompt.
\item \textbf{Code}: generate code conditioned on the clarified prompt, and
run the real, executable test suite associated with the problem.
\end{enumerate}

This constitutes round one. Round two repeats steps 1--4 once more,
conditioned on the round-one clarified prompt, to model an agent that gets
a second opportunity to resolve any remaining ambiguity before finalizing
its code. We report both round-1 and round-1+2 correctness throughout
Section~\ref{sec:results}, since, as we show, the value of a second
clarification round does not straightforwardly track the value of the
first.

\subsection{Design Rationale for the Probe-Host Model}

We use Qwen3-8B as the probe-host model throughout due to compute
constraints in this study; the ranking probe's forward-pass cost scales
with this model's size, not with the size of whichever model ultimately
writes the code. This decoupling is itself part of the approach's intended
value proposition: ambiguity localization can, in principle, be performed
by a small, cheap model regardless of which (possibly much larger and more
expensive) model is used for code generation, whereas a self-consistency
baseline built on the coding model's own repeated sampling ties uncertainty
estimation cost directly to that larger model's per-token cost. We evaluate
this decoupling directly in Section~\ref{sec:results} by pairing the same
ranking decisions with two different coding models of different capability.
